\chapter{电功与电能}\label{ch:g9:electric-power-energy}

电费账单既不收取伏特，
也不收取安培——它收取的是你家电器一整个月里
悄悄相乘积累下来的东西。本章将这两个
电学量合成第三个量，给快与慢的消耗
起上正式的名字，并最终落到
门口的电表上：这是可以与真实金钱对账的物理。

\section{功率：输送的快慢}

\begin{definition}[功率与瓦特]\label{def:g9:electric-power-energy:power}
器件的\emph{功率}\index{功率} $P$ 表示其转换能量的
快慢，即每秒转换的焦耳数。其单位为
\emph{瓦特}\index{瓦特}（\unit{W}）：每秒一焦耳，
这一单位是为纪念蒸汽时代的工程师而命名的；
\emph{千瓦}（\qty{1}{kW} $=$ \qty{1000}{W}）用于厨房里的大功率电器。
功率不是能量的多少，而是
能量消耗的\emph{快慢}——好比水龙头的流量，而非水桶里的水量。
%
\end{definition}

\begin{proposition}[电功率]\label{prop:g9:electric-power-energy:UI}
设器件两端的电压为 $U$，通过的电流为 $I$，
则器件获得的电功率为
\[
P = U \times I ,
\]
单位为瓦特，即伏特与安培之积。该式可
在三个方向上使用，其第二种用法用于保险丝盒的计算：
$I = P/U$ 给出额定功率所对应的
用电器从电网取用的电流。
\end{proposition}

\begin{example}[家用电器的功率阶梯]\label{ex:g9:electric-power-energy:ladder}
看一看家中各电器铭牌上的数值：指示灯，
\qty{1}{W}；现代灯，\qty{5}{W}（它的白炽前身
为获得同样亮度耗电六十瓦——差别全在于热量）；
笔记本电脑，\qty{50}{W}；电视机，\qty{100}{W}；
加热用水的洗衣机，\qty{2000}{W}；
电水壶，\qty{2300}{W}；即热式热水器，最高可达
\qty{9000}{W}。这把阶梯的规律是：凡以
\emph{加热}为主要用途的电器，功率都攀升到千瓦级——
保暖是最费能量的一种供能形式。
\end{example}

\begin{example}[保险丝盒的计算]\label{ex:g9:electric-power-energy:fuses}
接在 \qty{230}{V} 市电上的电水壶：$I = 2300 \div 230 =
\qty{10}{A}$——旧章中那台十安培的电器，
如今可以推导得出。\qty{9000}{W} 的热水器：
电流接近 \qty{40}{A}，需要专用的粗支路和断路器。另外，
火灾调查中那个超载的多插口插座：
将电水壶、热水器和电炸锅各自的 $P/U$ 电流相加，
超过支路的额定值——三个无辜的电器，
一桩算术罪行，正如指控所言。
\end{example}

\section{能量：输送的总量}

\begin{proposition}[由功率与时间求能量]\label{prop:g9:electric-power-energy:Pt}
功率为 $P$ 的器件工作时间为 $t$ 时，
转换的能量为
\[
E = P \times t ,
\]
单位为焦耳，即瓦特与秒之积。无论是快速放水
还是细水长流，水桶装水的多少都由这一乘积决定：一台 \qty{2300}{W} 电水壶
三分钟（$2300 \times 180 \approx \qty{4.1e5}{J}$），
消耗的能量多于一盏 \qty{5}{W} 灯一整天的（$5 \times
86400 = \qty{4.3e5}{J}$）——几乎恰好相等：
实际上，烧开一壶水恰好抵得上一天的好光照。
\end{proposition}

\begin{definition}[千瓦时]\label{def:g9:electric-power-energy:kwh}
家中的能量若以焦耳计数，数字显得笨拙，因此电表
改用\emph{千瓦时}\index{千瓦时}计量。
（\unit{kWh}）：功率为一千瓦的器件持续工作
一小时所转换的能量，
\[
\qty{1}{kWh} = 1000 \times 3600 = \qty{3.6e6}{J}
\]
——即三百六十万焦耳，正是账单上的每一个计费单位。
若 $P$ 以千瓦计、$t$ 以小时计，则 $E = P \times t$
可直接得到以千瓦时为单位的电能，
这正是电表所依据的公式。
\end{definition}

\begin{method}[电器耗电量的核算]\label{met:g9:electric-power-energy:audit}
对任一电器与任一月份：
\begin{enumerate}
  \item 从铭牌上读出功率 $P$（或计算
        $U \times I$);
  \item 据实估计其运行时间 $t$——先按每天，
        再折算为每月；
  \item 由 $E = P \times t$ 算出以千瓦时计的电能（$P$ 以
        \unit{kW}，$t$ 以小时计）；
  \item 乘以电价（约 $0.25$ 货币单位每
        \unit{kWh}），得出当月电费。
\end{enumerate}
每天用电水壶六分钟：$2.3 \times 0.1 \times 30 =
\qty{6.9}{kWh}$，即每月耗电。电视机每天三小时：
$0.1 \times 3 \times 30 = \qty{9}{kWh}$。小功率
长时间运行，可与大功率短时间运行的电器相当——
这正是核算中反复出现的一课。
\end{method}

\begin{omfigure}
\begin{tikzpicture}
\begin{axis}[omaxis, width=10.6cm, height=5.4cm,
    xlabel={时间（\unit{h}）}, ylabel={功率（\unit{kW}）},
    xtick={0,1,2,3}, ytick={0,1,2,2.3},
    ymin=0, ymax=2.9, xmin=0, xmax=3.4]
  \addplot[very thick, omThm] coordinates
    {(0,2.3) (0.1,2.3)};
  \addplot[omThm, fill=omThm, fill opacity=0.25, draw=none]
    coordinates {(0,0) (0,2.3) (0.1,2.3) (0.1,0)} \closedcycle;
  \addplot[very thick, omDef] coordinates {(0,0.1) (3,0.1)};
  \addplot[omDef, fill=omDef, fill opacity=0.25, draw=none]
    coordinates {(0,0) (0,0.1) (3,0.1) (3,0)} \closedcycle;
  \node[right, font=\footnotesize, omThm] at (axis cs:0.15,2.0)
    {电水壶：功率高，时间短};
  \node[above, font=\footnotesize, omDef] at (axis cs:1.9,0.14)
    {电视机：功率低，时间长};
\end{axis}
\end{tikzpicture}

{\small 能量即面积：功率乘以时间。水壶细高的窄条
与电视机晚间的长条可以围出相当的面积——
即相当的千瓦时。}
\end{omfigure}

\section{电表与电费账单}

\begin{example}[读懂账单]\label{ex:g9:electric-power-energy:bill}
门口的电表计量每一份进入房屋的
电能（千瓦时）；账单上的电费则是当月读数之差乘以电价。
一个普通家庭每月 $250$ 度用电，经拆分审计：
先是热水与暖气片（千瓦级
俱乐部），再是冷藏类电器——功率不大，但
始终在运行；然后是烹饪、洗衣、照明与
电子设备。审计的价值在于：找出真正的杠杆。
把十只旧的 \qty{60}{W} 灯泡换成 \qty{5}{W} 的，
节省的电能胜过把家中每个充电器
拔掉一千次所省下的电。
\end{example}

\begin{example}[待机“吸血鬼”]\label{ex:g9:electric-power-energy:standby}
小功率，长时间：电视机的待机指示灯
仅耗 \qty{0.5}{W}——但按整月 $720$ 小时计算：
$0.0005 \times 720 = \qty{0.36}{kWh}$。单个“吸血鬼”
无伤大雅；可若一户人家有二十个每个一瓦的耗电者，
每月白白耗费 $0.02 \times 720 \approx \qty{14}{kWh}$。
功率乘以时间，每一项都不能忽略：第二个因子
在第一个因子不起眼时，往往最容易被忽略。
\end{example}

\begin{remark}[功率在更广阔的世界中]\label{rem:g9:electric-power-energy:wider}
瓦特这一单位源于电学，度量着一切能量
的快慢。人体静息时的功率约为 \qty{100}{W}——
一只发光发热的“祖先灯泡”（想想拥挤教室的闷热：
三十名学生，相当于三千瓦的生物“热源”）。
骑车者可持续输出 \qty{200}{W}；汽车发动机可释放
几十千瓦；城镇的发电站，功率达数
亿瓦；落在小镇上的阳光，
远远超过这些——它是其他一切能量活动
所处的上限。下一章将追溯电网的瓦特
至它们的旋转之源。
\end{remark}

\section{练习}

\begin{exercise}[$\star$]\label{exo:g9:electric-power-energy:1}
定义功率及其单位——它像“水龙头”还是“水桶”？
写出本章的两个公式。
\end{exercise}

\begin{exercise}[$\star$]\label{exo:g9:electric-power-energy:2}
计算以下功率：一盏灯在 \qty{230}{V} 电压下
吸取 \qty{0.022}{A} 的电流；一台起动电动机在 \qty{12}{V} 电压下
吸取 \qty{100}{A} 的电流。
\end{exercise}

\begin{exercise}[$\star$]\label{exo:g9:electric-power-energy:3}
保险丝箱计算：一台 \qty{2300}{W} 的
电水壶与一台 \qty{700}{W} 的微波炉，在 \qty{230}{V} 下各吸取多大电流？
\end{exercise}

\begin{exercise}[$\star$]\label{exo:g9:electric-power-energy:4}
为什么加热类电器总是占据家用
用电“阶梯”的最高处？
\end{exercise}

\begin{exercise}[$\star$]\label{exo:g9:electric-power-energy:5}
换算：将 \qty{1}{kWh} 换成焦耳（写出两个换算因子）。
为什么电表的制造者弃用了焦耳？
\end{exercise}

\begin{exercise}[$\star$]\label{exo:g9:electric-power-energy:6}
能量核算：哪一个耗能更多，是 \qty{2}{kW} 的电暖器
工作半小时，还是 \qty{100}{W} 的电视机
连续开上整晚五小时？
\end{exercise}

\begin{exercise}[$\star$]\label{exo:g9:electric-power-energy:7}
审计一台机器：一台 \qty{50}{W} 的笔记本电脑，每天使用四小时，
共三十天，电价为每 \unit{kWh} $0.25$ 元。
求耗电量与电费。
\end{exercise}

\begin{exercise}[$\star$]\label{exo:g9:electric-power-energy:8}
门口的电表计量的是什么？账单如何把
电表读数转化为金额？
\end{exercise}

\begin{exercise}[$\star\star$]\label{exo:g9:electric-power-energy:9}
灯泡更换：十只 \qty{60}{W} 的旧灯泡与十只
\qty{5}{W} 的新灯泡相比，每晚点亮四小时，共三十晚。
计算两种灯泡各自的月耗电量及节省的电能，
以千瓦时计，电价为每千瓦时 $0.25$ 元。
\end{exercise}

\begin{exercise}[$\star\star$]\label{exo:g9:electric-power-energy:10}
一台吹风机的铭牌上写着“230 V，2000 W”。它的电流是多少？
它能否与 \qty{2300}{W} 的电水壶共用一条 \qty{16}{A} 的支路，
两者同时工作？请列出电流之和。
\end{exercise}

\begin{exercise}[$\star\star$]\label{exo:g9:electric-power-energy:11}
一辆电动汽车的电池储存 \qty{60}{kWh} 的电能，
这些电能能让 \qty{2.3}{kW} 的电水壶工作多久？
\qty{5}{W} 的灯泡又能亮多久？这两个结果说明了
交通运输的能量需求规模有多大？
\end{exercise}

\begin{exercise}[$\star\star\star$]\label{exo:g9:electric-power-energy:12}
电水壶的全程检验：将 \qty{1.5}{L} 的水从
\qty{20}{\celsius} 加热到沸腾，约需 \qty{5.0e5}{J}
的热量。为家用电水壶（\qty{2300}{W}）计时，
利用 $E = P t$ 预测其沸腾所需时间，再
解释为何实际的电水壶要稍慢一些——那些缺失的
焦耳去了哪里？请依据能量章节中最古老的定律作答。
\end{exercise}

\section{问题：家庭能量审计}

\begin{problem}[{周末问题 --- 家庭能量审计；
铭牌、电表与当月账单；找出
真正的杠杆}]\label{pb:g9:electric-power-energy:1}
凭借两个公式和电价（$0.25$ 元每
\unit{kWh}），全家审计一个月（30 天）。
清单：电水壶 \qty{2300}{W}，每天 $6$ 分钟；
电视机 \qty{100}{W}，每天 $4$ 小时；冰箱 ---
见下文特殊情形；十盏灯各 \qty{5}{W}，每天
$5$ 小时；洗衣机 \qty{2}{kW}，
隔天一次 $1$ 小时的热水循环；热水器 \qty{2.5}{kW}，每天
$2$ 小时；待机小耗电器共 \qty{10}{W}，始终开启。

\medskip\noindent\textbf{第一部分 --- 逐台计算。}
\begin{enumerate}
  \item 电水壶一个月的耗电量，单位为 \unit{kWh}。
  \item 电视机与十盏灯各自一个月的用电量。
  \item 洗衣机（$15$ 次循环）与
        热水器一个月的用电量。
  \item 待机小耗电器一个月的用电量 --- 这是
        起作用的第二个因素。
\end{enumerate}

\medskip\noindent\textbf{第二部分 --- 冰箱与
总计。}
冰箱压缩机运行时功率为 \qty{120}{W}，但
只占三分之一的时间，昼夜循环运行。
\begin{enumerate}[resume]
  \item 冰箱的等效平均功率，以及以
        \unit{kWh} 为单位的一个月用电量。
  \item 将审计表的七项相加。哪三项
        占主导地位？
  \item 按电价算出当月电费 --- 以及
        热水器所占部分的金额（以元计）。
  \item 月初电表读数为 $2150$ 度。
        预测月末的电表读数。
\end{enumerate}

\medskip\noindent\textbf{第三部分 --- 杠杆。}
\begin{enumerate}[resume]
  \item 方案一：缩短淋浴时间 --- 热水器
        每天的运行时间降为 \qty{1.5}{h}。每月节省的
        能量与电费各为多少？
  \item 方案二：公益宣传呼吁拔掉
        手机充电器（待机时每个 \qty{0.2}{W}）。
        审计一个充电器待机一个月的用电量，
        并与方案一对比。这一对比能说明
        审计与口号的区别是什么？
  \item 方案三：奶奶建议沿用旧习惯，
        即只按需要灌水 --- 水量减半，
        能量减半。电水壶一项的新用电量，以及
        一个月节省多少？
  \item 写出审计结论：三句话 ---
        说明千瓦时真正藏在何处，待机情形
        揭示了公式中的哪个因子，以及哪一项
        单一的家庭改变最划算。
\end{enumerate}
\end{problem}
